%% Appendices -- extended version only (\ifextended). ICRA does not accept a
%% separate supplementary PDF, so none of this may be load-bearing for review.

\section{Complete Pipeline Specification}
\label{app:spec}

\subsection{Coordinate conventions}
One metric, $z$-up world frame throughout: the ScanNet++ mesh/COLMAP
frame, floor near $z{=}0$. Camera poses are OpenCV world-to-camera
matrices parsed directly from \texttt{dslr/colmap/images.txt}, paired with
\texttt{PINHOLE} intrinsics from
\texttt{dslr/nerfstudio/transforms\_undistorted.json} and images from
\texttt{dslr/resized\_undistorted\_images}. The generator's canonical
object frame is $z$-up and roughly $[-0.5,0.5]^3$. A registration is
stored as a single $4{\times}4$ matrix $T$ with the isotropic scale
folded into the upper-left $3{\times}3$, so that
$x_{\text{world}} = sR\,x_{\text{canonical}} + t$.

A deliberate separation exists between the mesh used \emph{by} the
pipeline (raycasting, registration targets, collision background) and the
mesh used to \emph{evaluate} it. Ground-truth instance annotations are
indexed by vertex into the dataset mesh specifically, so the ablation of
Sec.~\ref{sec:ladder} replaces only the pipeline mesh; evaluation always
scores against the true dataset mesh. Conflating the two silently
corrupts every ground-truth comparison, and it is the single easiest way
to produce an ablation result that is quietly meaningless.

\subsection{Stage-by-stage parameters}

\begin{table}[htbp]
\caption{Instance discovery.}
\label{tab:s-discovery}
\centering\footnotesize\setlength{\tabcolsep}{4pt}
\begin{tabular}{@{}p{26mm}p{20mm}p{28mm}@{}}
\toprule
parameter & value & role \\
\midrule
frame stride & 12 & $\approx$28 frames on a typical scene \\
prompts & 16 (+14) & manipulable tier (+ furniture tier) \\
score threshold & 0.45 & SAM3 detection confidence \\
pixel stride & 4 & raycast grid inside each mask \\
cluster radius & $\max(0.5, 0.7 e_c)$\,m & about the median hit \\
min hits & 30 & before and after clustering \\
voxel size & 2\,cm & merge resolution \\
merge IoU & 0.25 & union-find over voxel sets \\
min frames seen & 2 & cross-view confirmation \\
min voxels & 25 & (half that, per proposal) \\
vertex snap & 3\,cm & points $\rightarrow$ mesh vertices \\
dedup overlap & 0.5 & of the \emph{smaller} vertex set \\
extent cap & $e_c$ (per category) & nominal maximum extent \\
\bottomrule
\end{tabular}\\[3pt]
{\footnotesize $e_c$ is the per-category nominal extent from the
vocabulary table (e.g.\ bottle 0.45\,m, keyboard 0.6\,m, cardboard box
1.2\,m); it defaults to 0.9\,m for unlisted labels.}
\end{table}

\textbf{Discovery} (Table~\ref{tab:s-discovery}). Proposals are formed
per (frame, prompt) pair; the full procedure is
Algorithm~\ref{alg:discovery}. The \emph{ground-truth enumeration} used in
benchmark mode is the drop-in substitute and emits an identical contract:
it filters dataset instances through a 54-alias whitelist, excludes
structural categories (wall, floor, ceiling, door, window, and 20 more),
and applies a size gate of $[0.03, e_c]$\,m with a global fallback
maximum of 0.8\,m.

\begin{table}[htbp]
\caption{Best-view selection.}
\label{tab:s-bestview}
\centering\footnotesize\setlength{\tabcolsep}{4pt}
\begin{tabular}{@{}p{26mm}p{20mm}p{28mm}@{}}
\toprule
parameter & value & role \\
\midrule
visibility samples & 120 & instance vertices sampled \\
min in-frame fraction & 0.6 / 0.3 & small / $>$1.2\,m objects \\
min visible fraction & 0.5 & first-hit agreement rate \\
raycast tolerance & 2\,cm & 4\,cm for TSDF meshes \\
min bbox side & 48\,px & rejects distant views \\
geometric score & $\mathrm{vis}^2\sqrt{wh}$ & visibility $\times$ size \\
sharpness re-rank & top 8 & variance of Laplacian \\
sharpness clip & $[0.5,1.5]$ & vs.\ reference variance 80 \\
mask pad & $0.15\max(w,h)+8$\,px & crop margin \\
SAM3 evidence IoU & 0.15 & replace projected mask \\
min mask pixels & 400 & viability floor \\
\bottomrule
\end{tabular}
\end{table}

\textbf{Best-view selection} (Table~\ref{tab:s-bestview}). The visibility
test casts one ray per sampled instance vertex from the candidate camera
and requires the returned depth to agree with the known point distance
within tolerance; this is what makes the score occlusion-aware rather
than merely a frustum test. The tolerance must be widened when the
pipeline mesh is a TSDF fusion rather than a laser scan, because vertex
positions are quantized to roughly the voxel scale and genuinely visible
points otherwise miss the window---this single parameter was responsible
for a large fraction of the yield gap in an early version of ablation
row~C, and is exposed as an environment override for that reason.

\textbf{Asset generation.} TRELLIS runs with a fixed seed of 42 and its
output is decimated to a 40k-triangle simulation mesh; the
full-resolution mesh and the generated Gaussians are retained separately,
since the simulation mesh is what CoACD and the URDF consume while the
Gaussians are what the compositing renderer uses. ReconViaGen consumes up
to 12 views selected from the visibility ranking with a minimum index gap
of 4 between chosen frames, which prevents the selection collapsing onto
a burst of near-identical frames from one dwell in the capture
trajectory.

\begin{table}[htbp]
\caption{Registration.}
\label{tab:s-align}
\centering\footnotesize\setlength{\tabcolsep}{4pt}
\begin{tabular}{@{}p{26mm}p{20mm}p{28mm}@{}}
\toprule
parameter & value & role \\
\midrule
asset samples & 20\,000 & surface points, source \\
target samples & 6\,000 & subsampled for speed \\
yaw grid & $10^\circ$ & 36 candidates \\
scale multipliers & 7 in $[0.5,1.25]$ & per yaw \\
extent percentiles & 2 / 98 & robust extent \\
$z$ init percentile & 2 & visible bottom \\
clip distance $c$ & 3\,cm & chamfer saturation \\
ICP max iterations & 60 & point-to-point \\
upright snap & $<15^\circ$ tilt & project to yaw-only \\
scale polish & 13 in $\pm15\%$ & 2 alternating rounds \\
scale clamp & $[0.6,1.4]\times s_{\text{init}}$ & vs.\ extent init \\
translation rounds & 3 & inlier-mean shifts \\
\bottomrule
\end{tabular}
\end{table}

\textbf{Registration} (Table~\ref{tab:s-align}, Algorithm~\ref{alg:align}).
Note that the F1 evaluation uses the full 20k target samples even though
alignment is run against a 6k subsample; the subsample is a speed
optimization for the $36\times7$ candidate sweep and must not leak into
scoring.

\begin{table}[htbp]
\caption{Removal and completion.}
\label{tab:s-removal}
\centering\footnotesize\setlength{\tabcolsep}{4pt}
\begin{tabular}{@{}p{26mm}p{20mm}p{28mm}@{}}
\toprule
parameter & value & role \\
\midrule
instance radius & 3\,cm & selector (i) \\
asset radius & 2.4\,cm & selector (ii), 5k asset samples \\
vote AABB dilation & 15\,cm lateral, & selector (iii) candidates \\
 & $-5$/$+15$\,cm in $z$ & \\
vote threshold & $\geq2$ views & opacity-agnostic \\
related views $K$ & 3 & top-ranked by visibility \\
plane ring & 10\,cm annulus & outside the footprint \\
plane $z$ band & 3\,cm & about the object's lowest vertex \\
MAD trim & $2.5\times$, $\leq2$ rounds & floor of 50 points \\
fill grid & 5\,mm & on the plane \\
hull offset & 3.5\,cm & exceeds the carve radius \\
disk scales & $4{\times}4{\times}0.8$\,mm & thin axis $\parallel$ normal \\
plane offset & 0.5\,mm & along $+n$, avoids z-fighting \\
refine iterations & 1500 & diff.\ rasterization \\
render scale & 0.5 & during refinement \\
initial opacity & $\sigma^{-1}(0.9)$ & logit 2.2 \\
inpaint crop pad & 0.35 & relative to mask bbox \\
inpaint crop max & 1024\,px & longest side \\
diffusion steps & 30 & Qwen-Image-Edit \\
mask dilation & 9\,px & SAM3 union masks \\
\bottomrule
\end{tabular}
\end{table}

\textbf{Removal and completion} (Table~\ref{tab:s-removal}). Colour
initialization projects each fill point into its object's primary
inpainted view and reads that pixel, converting to the zeroth
spherical-harmonic coefficient by $(\text{rgb}-0.5)/C_0$ with
$C_0 = 0.2820948$; the mean of the five nearest surviving Gaussians is
only a fallback for points that project outside the frame. Initializing
from neighbour colours alone left the optimization far from its target and
1500 iterations of inconsistent supervision could not recover it.

\textbf{Physics annotation.} CoACD runs with concavity threshold 0.05 and
a hull budget of 16, doubled to 32 for objects whose largest world
dimension exceeds 1.2\,m. The LLM is prompted with the category label, the
metric axis-aligned bounding box, and a fixed context (``as found on an
office desk''), and constrained to emit a single JSON object with keys
\texttt{mass\_kg}, \texttt{friction}, \texttt{restitution}. Sampling is
disabled (greedy decoding) for reproducibility. A response is accepted
only if mass $\in[0.005,20]$\,kg, friction $\in[0.2,1.0]$ and restitution
$\in[0,0.5]$; otherwise a per-category density-table fallback is used, and
the record is tagged with which path produced it so the fallback rate is
auditable. Inertia is the uniform-density box tensor of the world
dimensions and the inertial origin is the canonical mesh centroid.

\textbf{Simulation and export.} PyBullet runs at 240\,Hz. Settling uses
velocity-zeroed stepping (velocities reset to zero after every step) for
2\,s, which resolves initial interpenetration without launching objects.
The background is the carved scan mesh decimated to 300k triangles. MJCF
export emits each CoACD part as a separate mesh geom---MuJoCo convexifies
each geom independently, so a convex decomposition survives the export
exactly---plus a per-object static support slab of 2\,cm thickness. The
floor plane is placed 5\,cm below the lowest object bottom rather than at
$z{=}0$, which is only correct for scenes whose floor is the world origin.
Degenerate hulls with fewer than four vertices are dropped before
emission, since MuJoCo's compiler rejects them.

%======================================================================
\section{Metric Definitions}
\label{app:metrics}

\subsection{Discovery F1}
Let $\mathcal{P}$ be predicted instances and $\mathcal{G}$ ground-truth
instances, each a set of mesh vertex indices. For a threshold $\theta$,
build the bipartite graph on pairs with
$\mathrm{IoU}(p,g) = |p\cap g|/|p\cup g| \geq \theta$ and greedily match
in descending IoU, one-to-one. With $M$ matches,
$\mathrm{precision} = M/|\mathcal{P}|$,
$\mathrm{recall} = M/|\mathcal{G}|$, and F1 is their harmonic mean. We
report $\theta \in \{0.25, 0.5\}$. Ground truth is restricted to
whitelist-sized instances so that predictions are not penalized for
missing categories outside the prompt set.

\subsection{Geometry F1@$\tau$}
Sample $N=20{,}000$ points on the registered asset surface ($A$) and on
the reference surface ($R$), the latter being the ground-truth instance
submesh (triangles all of whose vertices belong to the instance):
\begin{align}
\mathrm{prec} &= \tfrac{1}{|A|}\textstyle\sum_{a \in A}
   \mathbb{1}\!\left[\min_{r\in R}\|a-r\| < \tau\right],\\
\mathrm{rec}  &= \tfrac{1}{|R|}\textstyle\sum_{r \in R}
   \mathbb{1}\!\left[\min_{a\in A}\|r-a\| < \tau\right],
\end{align}
with $\text{F1} = 2\,\mathrm{prec}\cdot\mathrm{rec}/(\mathrm{prec}+\mathrm{rec})$
at $\tau\in\{20,40\}$\,mm. We also record the symmetric mean chamfer
$\tfrac{1}{2}(\overline{d}_{A\to R} + \overline{d}_{R\to A})$.

\textbf{Pooling.} All headline geometry numbers are pooled over objects:
every tier-A/B object contributes one F1 value to one mean. The per-scene
macro mean reads 0.636 against the pooled 0.708 on the same data, because
five scenes yield zero objects and contribute a hard zero.

\textbf{Reference choice.} For self-discovered instances we report two
scorings. \emph{Matched-GT} first matches each discovered instance to a
ground-truth instance at vertex IoU $\geq0.25$ and scores the registered
asset against that ground-truth submesh; unmatched instances are excluded
from the F1 but still counted in yield. \emph{Own-extraction} scores
against the point set the pipeline itself extracted for that instance.
Only the former is an independent measurement, and
Sec.~\ref{sec:ladder} gives a case where the two disagree in sign.

\subsection{Stability ratio}
Following PhyRecon~\cite{phyrecon}, SR is the fraction of non-rejected
assets that survive an isolated drop test. Concretely: reset the
simulation, gravity $-9.81\,\text{m/s}^2$, timestep $1/240$\,s, add an
infinite plane, spawn the object's URDF at $z=5$\,mm with inertia read
from the file; step $2\times240$ times with linear and angular velocity
reset to zero after each step (settling); record the pose $p_0$; step a
further $2\times240$ times freely; record $p_1$. The object is stable iff
$\|p_1-p_0\| < 3$\,cm and $p_{1,z} > -5$\,cm.

The pose is read in the \emph{canonical link frame}, not PyBullet's
inertial/COM frame. PyBullet reports base poses in the COM frame, and for
an object whose centroid is offset from its link origin the two differ by
a fixed vector that rotates with the body, so any tilt registers as
spurious translation. Measured on our benchmark, this convention
reclassifies 3 of 457 objects (GT-driven) and 4 of 671 (automatic) --- it
must be stated, but it is not large. The larger part of the difference
between our current 77.2\% and the 79.9\% an earlier draft reported is the
\emph{asset population}: holding the frame convention fixed, the
residual-gated hybrid moves stability from 79.9\% to 76.6\%, and the
frame fix then returns $+0.7$ points. We give the decomposition because
it is the kind of change that is otherwise invisible between paper
versions.

\subsection{Appearance metrics}
Frames are the intersection of each scene's official
\texttt{dslr/train\_test\_lists.json} \emph{test} list with
COLMAP-registered frames, subsampled to at most 8 by uniform index
spacing. PSNR is computed on RGB in $[0,1]$, SSIM via \texttt{skimage}
with \texttt{channel\_axis=2} and \texttt{data\_range=1.0}, and LPIPS with
the AlexNet backbone on inputs mapped to $[-1,1]$. Four variants are
rendered per frame: background splat alone, composite twin (background
plus every tier-A/B asset's Gaussians transformed by its registration),
and---when a cleaned background exists---clean background alone and
assets composited onto it. Asset Gaussians are spherical-harmonic-padded
to the background's degree before concatenation.

An earlier version of this evaluation selected held-out frames by uniform
spacing over \emph{all} registered frames, which mostly selected frames
the splat had trained on and inflated every PSNR. The current protocol
falls back to that behaviour only when a scene ships no usable test list,
and prints a warning identifying the scene when it does.

\subsection{Verification battery}
All four checks run on renders of the cleaned background, at half
resolution, for both supervision views and two held-out neighbours offset
$\pm7$ frames along the sorted capture trajectory.

\emph{Alpha coverage} is the mean rendered opacity inside the object's 2D
mask. Values near 1 mean the region is filled; low values mean the removal
punched a hole through to the background.

\emph{Support residual} intersects each masked pixel's viewing ray with
the fitted support plane, and reports the median absolute difference
between the rendered expected depth and the plane depth, over pixels where
both exceed 0.1\,m.

\emph{Residual re-detection} runs SAM3 with the object's own text prompt on
both the original and cleaned renders. A detection counts if it overlaps
the scoring region by more than 30\% of the detection's own area, and the
score is the maximum confidence among counting detections. Removal passes
if the maximum post-removal score across views is below 0.5. On
supervision views the scoring region is the object mask; on held-out views
it is the 2D bounding box of the removed Gaussian centers projected into
that view, dilated by 15\% of the image diagonal. Held views for which
that region cannot be constructed (fewer than 20 projected centers in
frame) are excluded rather than scored whole-frame.

%======================================================================
\section{Registration Stress Suite}
\label{app:stress}
The suite (shipped as a runnable test) draws 6 shapes $\times$ 3 partial
views under two conditions---clean, and corrupted with 2\,mm Gaussian
noise plus a 5\% outlier fraction simulating depth bleed---and scores
recovery against tolerances of 5\% scale, 5\,mm translation and
$10^\circ$ yaw. Partial views are generated by sampling 40k points on the
complete shape and culling to those visible from a random camera, then
subsampling to 6\,000 points to match the registration target size used
by the benchmark.

Prismatic and feature-bearing shapes (box, L-shape, mug) recover
reliably. Cones and capsules underestimate scale by 7--50\%. The suite
establishes that this is a yaw-grid initialization local minimum rather
than data ambiguity, because the ground-truth pose scores strictly
\emph{better} than the recovered pose under the same objective---the
optimizer is not being misled about which pose is correct, it never visits
it. Contamination is also not the driver: clean and corrupted conditions
fail on nearly identical cases, which is independent evidence that the
clipped objective is doing its job. Crucially there are zero $180^\circ$
flips in either condition; that is the failure mode a one-way chamfer
exhibits, and it is the reason the objective is symmetric.

%======================================================================
\section{BEHAVIOR-1K Coverage Protocol}
\label{app:behavior}
The feasibility check of Sec.~\ref{sec:behavior1k} proceeds per activity.
We parse each BDDL problem's object list and keep categories whose
BEHAVIOR synset is rigid-body-like, excluding architectural and fixture
categories that any indoor scan already contains (walls, floors, doors,
windows, and plumbed or wired fixtures such as sinks, stoves and
cabinetry)---these are present in a twin whether or not we generate an
asset for them, so counting them would inflate coverage. Matching against
our vocabulary is deliberately strict: a compound synset such as
\emph{detergent\_\_bottle} counts as covered only if the \emph{contents},
not merely the container, correspond to a label we handle; a plain
\emph{bottle} does not satisfy it. We report two coverage levels: against
the full prompt/whitelist vocabulary, and against the strictly smaller set
of labels the pipeline actually produced across the 50 validation scenes.
Room plausibility is scored separately against a list of room types a scan
dataset like ScanNet++ would plausibly contain, and against an explicit
implausible list (garden, garage, and similar). An activity is ``fully
groundable'' only if every one of its required categories is in the
discovered set.

%======================================================================
\section{Environment and Reproduction}
\label{app:env}

\subsection{Three interpreters}
The pipeline requires three python environments with mutually incompatible
builds, and the launcher resolves all three:
\begin{itemize}
\item \textbf{main} (torch 2.4.1+cu124): everything except the two below.
\item \textbf{segmentation} (torch 2.10): SAM3, and Qwen-Image-Edit, which
  calls \texttt{enable\_gqa} and therefore needs torch $\geq2.5$ plus a
  newer \texttt{diffusers} than the main environment tolerates. On a
  48\,GB card it additionally requires sequential CPU offload.
\item \textbf{rasterization} (python 3.10): every gsplat CUDA render. The
  only prebuilt gsplat wheel is \texttt{cp310}, and a JIT build for the
  main environment's python 3.11 is not available on our nodes.
\end{itemize}
A fourth, isolated environment pins the MaskClustering baseline's
dependencies. The segmentation environment carries numpy 2.x, where
\texttt{ndarray.ptp} was removed, so code shared between environments must
use \texttt{np.ptp}.

\subsection{Running}
Each stage is a module. The four entry points are a GT-driven benchmark
run, a fully automatic run, a single-frame zero-shot run, and the
removal/inpainting pass; each is a shell script that sources a shared
environment file and calls stages in order with resumable guards, so an
interrupted scene restarts at the first incomplete stage. Scene selection,
output directory, discovery mode, mesh source and vocabulary tier are all
environment variables. Fleet runs are slurm array jobs that stride over
the validation split and skip scenes whose completion markers exist.

Every table in this paper that reports measured numbers is generated
directly from the outputs tree by a single script, so the paper cannot
drift from the results it describes.

\subsection{Timing methodology}
Reported costs are the sum of every GPU-second spent producing the final
artifacts, including partial runs that were later resumed, rather than the
wall-clock of a single clean pass. This is the conservative convention: it
charges us for restarts. Per-stage timings are appended to a per-scene
timings file by the launcher itself, so the accounting does not depend on
post-hoc estimation. We do not report H200 numbers: the batch partition
hosting those cards does not mount the filesystem this project lives on,
so we could not run the pipeline there at all.

%======================================================================
\section{Baseline Reimplementation Details}
\label{app:baselines}

\subsection{MaskClustering}
We run the released implementation with its published ScanNet++
configuration, changing only the mask source. Our harness exports SAM3
masks in the format its dataset loader expects, converts its predicted
object dictionary back into our vertex-set contract, and scores it with
the identical discovery-F1 code used for our own instances. Two
confounds are worth stating. First, its protocol ingests every second
DSLR frame while ours strides every twelfth; we therefore report it at
both its native density and at our matched 28-frame budget, and the
matched number is the fair one. Second, its hyperparameters were tuned
for Cropformer masks and we did not re-tune them for SAM3; a tuned
version would likely close some of the remaining gap at IoU\,0.5.

\subsection{FlashSplat}
We reimplemented FlashSplat's optimal mask-to-Gaussian assignment against
our own Gaussian representation, so that both methods select from
precisely the same Gaussian set. Its formulation assumes a view set that
covers the object; our per-object view sets are small and object-centric,
which leaves many Gaussians with no evidence in any view. Assigning those
to the foreground produced obviously wrong removals, so zero-evidence
Gaussians are assigned to the background. The comparison in the main paper
is against our \emph{final} removal set including the mask vote; an
earlier draft compared against a pre-vote reference and reported a lower
agreement.

\subsection{Inpainting backends}
Qwen-Image-Edit is run as the Edit-Plus pipeline rather than the
inpainting pipeline, which applies a different and, for this task,
incorrect conditioning. LaMa serves as the fallback and runs on CPU to
avoid contending with resident diffusion weights. Backends are recorded
per object and per view so that the comparison in the main paper is
computed over exactly matched frames.

%======================================================================
\section{Failure Analysis}
\label{app:failures}

\textbf{Transparent objects} are the dominant failure mode and appear at
three separate stages. Their scan geometry is incomplete, so the projected
mask is wrong (mitigated by SAM3 image evidence); their splat
representation is a diffuse low-opacity cloud detached from the surface,
so geometric removal selectors miss them (mitigated by the multi-view
vote); and multi-view generation collapses on them outright, which is what
the registration-residual gate is for. After all three mitigations they
remain the two residual re-detection failures on the pilot scene, and the
category sits at 51\% yield.

\textbf{Box-scale holes.} The support-plane fill is designed for objects
whose footprint is small relative to the supporting surface. Cardboard
boxes violate this: their removal leaves a hole comparable in size to the
visible desk, and the plane fit's surrounding annulus can fall off the
support entirely. This is the same scale limit that keeps furniture out of
the main protocol, and it shows up in the verification battery as alpha
coverage of 0.83--0.86 and support residuals up to 0.8\,m.

\textbf{Deformables treated as rigid.} Bags (F1 0.525) and backpacks
(0.433) are the worst geometry categories despite high yield. A rigid
asset cannot represent an object whose shape in the scan is one draped
configuration among many, and the registration residual has no way to
express that the shape itself is wrong rather than mis-posed.

\textbf{Occlusion-limited categories.} Books (21\% yield) and jars (18\%)
fail at the opposite end: the ones that survive score well (0.809, 0.711),
but most never obtain a usable best view. The library scene alone
contributes roughly 200 shelved books, 189 correctly rejected; that scene
also requires roughly three hours and exceeded the wall clock of our
standard job class.

\textbf{Support modelling in MJCF.} Because the MJCF export approximates
the background with per-object slabs rather than the scan mesh, an object
that drifts a few centimetres can slide off its own slab and fall further
than it would on a continuous desk. This inflates apparent instability in
the MuJoCo export relative to the PyBullet export, which collides against
the carved scan. We report stability from the isolated plane drop test, so
this affects only the qualitative MuJoCo demonstrations.

%======================================================================
\section{Per-Scene Results}
\label{app:perscene}
Table~\ref{tab:perscene} lists every validation scene: instance count,
tier A and B counts, pooled geometry F1@20\,mm over tier-A+B objects, the
corrected link-frame drop test, background and composite-twin PSNR on the
official test split, and wall-clock minutes. The spread is wide---per-scene
F1 ranges over roughly 0.4 and per-scene PSNR over roughly 10\,dB---which
is why the main paper pools per object rather than averaging per scene.

\begin{table*}[htbp]
\caption{Per-scene results, GT-driven mode, all 50 ScanNet++ validation
scenes. \emph{stab.} is the corrected link-frame drop test (stable/tested);
\emph{bg/twin} is background and composite PSNR on the official test split.}
\label{tab:perscene}
\centering\scriptsize
% generated by simany.eval.make_paper_tables -- do not edit
\begin{tabular}{lrrrrrrr}
\toprule
scene & $n$ & A & B & F1@20 & stab. & bg/twin & min \\
\midrule
\texttt{09c1414f1b} & 47 & 24 & 3 & 0.652 & 23/27 & 24.3/24.2 & 33 \\
\texttt{0d2ee665be} & 24 & 18 & 1 & 0.706 & 15/19 & 32.0/29.0 & 19 \\
\texttt{13c3e046d7} & 0 & 0 & 0 & 0.000 & -- & 29.6/29.6 & 2 \\
\texttt{1ada7a0617} & 11 & 9 & 1 & 0.877 & 10/10 & 29.8/29.4 & 9 \\
\texttt{21d970d8de} & 3 & 2 & 0 & 0.724 & 2/2 & 30.7/30.7 & 4 \\
\texttt{25f3b7a318} & 29 & 15 & 5 & 0.617 & 16/20 & 28.7/27.0 & 21 \\
\texttt{27dd4da69e} & 17 & 11 & 1 & 0.709 & 10/12 & 34.9/34.4 & 14 \\
\texttt{286b55a2bf} & 10 & 3 & 1 & 0.595 & 3/4 & 32.9/32.4 & 6 \\
\texttt{31a2c91c43} & 0 & 0 & 0 & 0.000 & -- & 33.3/33.3 & 1 \\
\texttt{3864514494} & 5 & 3 & 0 & 0.790 & 2/3 & 29.6/29.3 & 4 \\
\texttt{38d58a7a31} & 2 & 2 & 0 & 0.665 & 2/2 & 27.1/27.1 & 3 \\
\texttt{3db0a1c8f3} & 49 & 26 & 6 & 0.645 & 25/32 & 28.4/28.2 & 34 \\
\texttt{3e8bba0176} & 6 & 5 & 0 & 0.885 & 5/5 & 29.3/28.8 & 5 \\
\texttt{3f15a9266d} & 216 & 21 & 2 & 0.773 & 13/23 & 26.8/26.6 & 90 \\
\texttt{40aec5fffa} & 2 & 1 & 1 & 0.429 & 0/2 & 32.4/32.2 & 3 \\
\texttt{45b0dac5e3} & 7 & 6 & 0 & 0.887 & 3/6 & 33.2/33.0 & 6 \\
\texttt{5748ce6f01} & 1 & 1 & 0 & 0.975 & 1/1 & 24.4/24.4 & 3 \\
\texttt{578511c8a9} & 25 & 21 & 0 & 0.854 & 18/21 & 26.3/26.1 & 19 \\
\texttt{5942004064} & 16 & 13 & 0 & 0.750 & 11/13 & 28.3/28.1 & 16 \\
\texttt{5eb31827b7} & 8 & 5 & 1 & 0.814 & 3/6 & 27.7/27.3 & 6 \\
\texttt{5ee7c22ba0} & 9 & 3 & 2 & 0.497 & 5/5 & 31.5/30.5 & 6 \\
\texttt{5f99900f09} & 17 & 11 & 5 & 0.559 & 13/16 & 24.7/24.5 & 16 \\
\texttt{6115eddb86} & 3 & 1 & 0 & 0.788 & 1/1 & 28.0/27.7 & 3 \\
\texttt{7831862f02} & 5 & 4 & 1 & 0.489 & 4/5 & 28.8/28.1 & 7 \\
\texttt{7b6477cb95} & 18 & 12 & 1 & 0.805 & 11/13 & 31.5/30.4 & 12 \\
\bottomrule
\end{tabular}
\hfill
\begin{tabular}{lrrrrrrr}
\toprule
scene & $n$ & A & B & F1@20 & stab. & bg/twin & min \\
\midrule
\texttt{7bc286c1b6} & 1 & 0 & 0 & 0.000 & -- & 34.6/34.6 & 1 \\
\texttt{825d228aec} & 9 & 5 & 3 & 0.580 & 5/8 & 31.4/30.6 & 7 \\
\texttt{9071e139d9} & 33 & 29 & 0 & 0.838 & 24/29 & 23.6/23.5 & 26 \\
\texttt{99fa5c25e1} & 7 & 6 & 0 & 0.610 & 5/6 & 21.8/21.6 & 6 \\
\texttt{a24f64f7fb} & 5 & 2 & 0 & 0.754 & 1/2 & 24.8/24.4 & 4 \\
\texttt{a8bf42d646} & 1 & 1 & 0 & 0.595 & 1/1 & 26.4/26.4 & 3 \\
\texttt{a980334473} & 4 & 2 & 1 & 0.501 & 3/3 & 32.7/28.8 & 4 \\
\texttt{ac48a9b736} & 9 & 8 & 1 & 0.694 & 9/9 & 27.4/27.3 & 7 \\
\texttt{acd95847c5} & 24 & 14 & 6 & 0.615 & 14/20 & 26.2/24.5 & 18 \\
\texttt{b0a08200c9} & 5 & 5 & 0 & 0.969 & 5/5 & 32.7/32.2 & 5 \\
\texttt{bcd2436daf} & 3 & 1 & 0 & 0.790 & 1/1 & 30.5/30.5 & 3 \\
\texttt{bde1e479ad} & 1 & 1 & 0 & 0.769 & 1/1 & 28.9/28.9 & 3 \\
\texttt{c49a8c6cff} & 14 & 11 & 0 & 0.769 & 9/11 & 26.8/26.5 & 10 \\
\texttt{c4c04e6d6c} & 2 & 1 & 0 & 0.629 & 1/1 & 28.9/28.8 & 3 \\
\texttt{c50d2d1d42} & 18 & 16 & 0 & 0.764 & 12/16 & 30.4/29.6 & 14 \\
\texttt{c5439f4607} & 7 & 5 & 1 & 0.716 & 5/6 & 29.1/28.6 & 7 \\
\texttt{cc5237fd77} & 12 & 8 & 1 & 0.759 & 9/9 & 23.0/22.9 & 7 \\
\texttt{d755b3d9d8} & 32 & 24 & 3 & 0.732 & 21/27 & 26.8/26.2 & 32 \\
\texttt{e398684d27} & 24 & 14 & 4 & 0.535 & 13/18 & 32.6/30.2 & 29 \\
\texttt{e7af285f7d} & 1 & 0 & 0 & 0.000 & -- & 26.2/26.2 & 3 \\
\texttt{f2dc06b1d2} & 6 & 2 & 3 & 0.446 & 4/5 & 29.5/29.1 & 6 \\
\texttt{f3685d06a9} & 0 & 0 & 0 & 0.000 & -- & 35.1/35.1 & 1 \\
\texttt{f3d64c30f8} & 26 & 16 & 3 & 0.698 & 14/19 & 28.9/25.5 & 23 \\
\texttt{f9f95681fd} & 5 & 4 & 0 & 0.742 & 4/4 & 26.0/25.6 & 8 \\
\texttt{fb5a96b1a2} & 10 & 7 & 1 & 0.791 & 8/8 & 25.7/25.6 & 8 \\
\bottomrule
\end{tabular}

\end{table*}

%======================================================================
\section{Engineering Findings}
\label{app:engineering}
These cost real time and are not recoverable from the method description.

\emph{Ordering.} CoACD must be imported before torch; the reverse order
segfaults through a bundled OpenMP clash.

\emph{Frames.} PyBullet reports base poses in the COM frame while URDF
authors think in the link frame. Every drift measurement must convert.

\emph{Composite targets.} When several objects are removed from one frame,
the refinement target must be composited by pasting each object's
inpainted pixels inside its own mask with a fixed priority. A
last-writer-wins composite bakes an un-removed neighbour's appearance into
the overlap band, because each single-object inpaint still depicts the
other object.

\emph{Supervision views.} Supervising the fill on every related view
produces a blotchy mean, because independent diffusion inpaints of
different frames invent mutually inconsistent texture. Supervising only on
primary views, and letting geometry propagate, is markedly better.

\emph{Caches.} Any stage that snapshots its inputs behind a one-time
marker will happily serve a stale snapshot after an upstream fix. The
hybrid generation stage has such a marker, and a rerun after a
registration-frame fix silently reproduced the old verdicts until the
snapshot directories were cleared.

\emph{Frame conventions.} Alignment code written for a $z$-up world
silently mismatches which axis is ``up'' when handed camera-frame data,
and the failure is not loud: it manifests as a size-sanity rejection with
a plausible-looking inflated scale, not as a crash.

\emph{Evaluation splits.} Selecting ``held-out'' frames by uniform spacing
over all registered frames silently selects training views. Always
intersect with the dataset's own declared test list, and fail loudly when
it is absent.
